\chapter{Related Work}

This chapter reviews the literature most relevant to the thesis. It does not aim to survey medical imaging broadly, but focuses on two areas that frame the present work. The first is CT-based mucus plug assessment, including clinical scoring and recent automated analysis. The second is methodological work on weakly supervised CT learning, multiple-instance learning, and transfer learning in limited-data settings. The chapter ends by identifying the gap addressed by this thesis, a patient-level mucus-burden prediction from chest CT when only coarse patient-level labels are available.


\section{CT-Based Mucus Plug Assessment}

\subsection{Clinical CT Scoring of Mucus Plug Burden}
The literature most closely aligned with this thesis is the CT-based mucus-plug scoring literature in asthma. Dunican et al.\ established a segment-based visual CT score and showed that greater mucus plug burden was associated with worse airflow obstruction and eosinophilic or type 2 airway inflammation \cite{dunican2018mucus}. This work is important here because it shows that mucus plug burden can be assessed systematically from CT and interpreted at the patient level, rather than only as a localized imaging finding. The related review by Dunican, Watchorn, and Fahy places mucus plugging in a broader clinical context by linking it to airway obstruction and disease mechanisms in asthma \cite{dunican2018autopsy}. Together, these studies frame mucus plugging as a clinically relevant burden measure that can be assessed from chest CT.

For this thesis, the key point is that the target is not the presence or absence of one visible plug, but the aggregate burden of plugging across the lungs. The original clinical scoring literature therefore justifies the patient-level endpoint used in this work. What it does not address is how such a score might be estimated automatically when the image evidence is distributed across a volumetric scan and only coarse labels are available.


\subsection{Recent Clinical and Pathophysiologic Work on Mucus Plugs}
More recent clinical literature suggests that mucus plugs remain an active pathophysiologic and imaging-relevant research theme beyond the original CT scoring studies. Hsieh et al.\ report that mucus plugs correlate with small airway remodelling in asthma, which strengthens the broader interpretation of mucus plugging as a marker of meaningful airway abnormality rather than only a descriptive CT finding \cite{hsieh2025smallairwaymucus}. This is important because it reinforces the idea that mucus-plug burden reflects underlying airway pathology that is relevant to disease severity and structural change.

From the perspective of the present thesis, such work helps broaden the significance of the target. The thesis does not aim to contribute directly to airway biology or histopathologic interpretation, but it benefits from a literature that continues to show that mucus plugs are clinically meaningful and mechanistically connected to disease processes. This matters because a computational prediction problem is more compelling when the quantity being predicted is already supported by a broader clinical and pathophysiologic literature.


\subsection{Recent Computational Work on Automated Mucus Plug Analysis}
Recent work has also begun to address automated computational analysis of mucus plugs on chest CT more directly. van der Veer et al.\ study automatic AI-based quantification of airway-occlusive mucus plugs in COPD and relate that quantification to all-cause mortality \cite{vanderveer2025aicopdmucus}. This is relevant because it demonstrates that automated quantification of mucus-plug burden is not only technically feasible, but can also be tied to clinically important outcome analysis. The paper therefore strengthens the broader case that automated mucus-plug analysis can have real clinical relevance rather than being only a technical image-processing exercise.

Raut et al.\ provide a different kind of contribution by validating an AI-based automated PRAGMA and mucus plugging algorithm in pediatric cystic fibrosis \cite{raut2025pragmaai}. Here the relevance lies in showing that automated mucus-plug analysis is being pursued across disease settings and that recent work is increasingly concerned with reproducible, validated computational pipelines rather than only exploratory proof-of-concept models. Although pediatric cystic fibrosis differs substantially from the present thesis setting, the paper still supports the broader point that automated CT-based mucus-plug analysis is becoming an active area of applied research.

Pu et al.\ move even further toward automated lesion-level analysis by presenting a simulation-driven, annotation-free deep learning framework for automated detection and segmentation of airway mucus plugs on thin-slice non-contrast chest CT \cite{pu2026simulation}. Their approach is notable because it reduces dependence on voxel-wise manual training annotations by generating synthetic plug examples within segmented airways. This is methodologically interesting for the present thesis because it shows another way in which the field is trying to overcome annotation bottlenecks in mucus-plug analysis. At the same time, the task formulation remains different from that of the present work, Pu et al.\ address plug-level detection and segmentation, not patient-level prediction of mucus burden from weak labels.

Together these newer studies show that automated mucus-plug analysis from chest CT is now an active recent direction. They span several related but distinct computational goals: automated quantification, disease-specific validation, and lesion-level detection or segmentation. This is an important development in the literature because it indicates that the field is beginning to move beyond purely visual or manually scored assessments toward more reproducible computational pipelines. At the same time, the methodological diversity of these studies also highlights that there is not yet one settled computational formulation of the problem.


\subsection{Limits of Existing CT Mucus-Plug Literature for the Present Task}
Despite this progress, the scope of the existing CT mucus-plug literature still differs from the scope of the present thesis. The original clinical scoring papers are primarily interpretive, they define and apply a patient-level burden score and evaluate its physiological relevance. More recent computational work extends the literature toward automated quantification, validation, detection, and segmentation, but it still does not primarily address automated patient-level regression from chest CT under weak supervision.

This distinction matters because the present thesis does not attempt to replace the clinical literature with a new scoring definition. Rather, it takes an existing clinically meaningful burden concept as its starting point and asks whether a score of this general type can be estimated computationally from CT using a constrained machine learning pipeline. The gap is therefore not in the existence of CT-based mucus-plug scoring itself, but in the limited literature connecting that scoring logic to weakly supervised patient-level prediction.

From the perspective of this thesis, the existing literature establishes the relevance of the target but does not resolve how a predictive model should handle volumetric CT data when only coarse patient-level labels are available. Clinical scoring papers justify the endpoint, newer automated studies show that computational mucus-plug analysis is feasible, but neither group of studies directly addresses the specific setting considered here, which is a patient-level mucus-burden prediction from chest CT under weak supervision and limited data. The present work is positioned in that methodological gap.



\section{Learning Approaches Relevant to CT-Based Patient-Level Assessment}


\subsection{Weakly Supervised and Multiple-Instance CT Learning}
Weakly supervised medical-imaging studies typically address the problem that clinically meaningful labels are available at a coarser level than the image evidence itself \cite{ren2023weakly}. In such settings, the methodological challenge is not only to learn from limited supervision, but also to determine how local image evidence should be related to a case-level target when the relevance of individual slices, patches, or regions is not known in advance. This has made multiple-instance and bag-level formulations particularly relevant in medical imaging, because they provide a natural way to represent an examination as a collection of local instances while keeping prediction defined at the level of the whole case.

In CT, this often leads to formulations in which a volume is represented as a collection of slices or patches and a case-level decision is obtained from pooled instance-level information. DR-MIL, for example used a bag-of-slices formulation with deep feature representation and bag-level classification \cite{qi2021drmil}. Xue et al.\ used attention-based multiple-instance learning for CT-based COPD identification and emphasized that disease evidence may be spatially heterogeneous rather than uniformly visible across the volume \cite{xue2023copdmil}. These studies are relevant because they illustrate the methodological design space surrounding the present thesis, coarse labels, unequal instance informativeness, and explicit case-level pooling.

Recent chest-CT work indicates that this remains an active methodological direction. Djahnine et al.\ propose a weakly supervised approach for pathology detection and localization in 3D chest CT, showing that coarse scan-level supervision can still support meaningful abnormality identification and localization in volumetric data \cite{djahnine2024weakly}. Almakky and Yaqub likewise propose a weakly supervised and explainable framework for infection severity classification from volumetric chest CT without relying on dense segmentation labels \cite{almakky2025weakly}. Together, these studies reinforce the broader computational setting of chest CT learning under coarse supervision and heterogeneous volumetric evidence. At the same time, they remain methodologically different from the present thesis, Djahnine et al.\ focus on pathology detection and localization, and Almakky and Yaqub study explainable severity classification, whereas the present work addresses patient-level regression of mucus burden.

The cited studies also differ from the present thesis in an important respect, they address case-level classification or detection-oriented tasks rather than patient-level regression. Their outputs are disease labels, abnormality localization signals, or severity categories rather than continuous burden scores, and their clinical targets differ substantially from mucus plug grading. Their value for the present work is therefore comparative rather than directly transferable. They clarify which weakly supervised aggregation strategies are plausible when a volumetric CT examination must be summarized into one prediction, but they do not directly solve the task considered in this thesis.


\subsection{Transfer Learning in Limited-Data CT Settings}
A second methodological theme relevant to the present thesis concerns transfer learning under limited data. In many medical-imaging settings, the available target cohort is too small to support robust end-to-end feature learning from scratch, which has made pretrained backbones an important practical and methodological tool \cite{kim2022transfer}. This is particularly relevant here because the thesis operates under a small-$N$ constraint and studies whether compact pretrained feature extractors can support patient-level prediction from chest CT under weak supervision.

This literature does not directly solve the present task either. Most transfer-learning studies in medical imaging focus on classification rather than regression, and many operate in larger datasets or with target definitions that differ from patient-level mucus-burden scoring. Nevertheless, the transfer-learning literature is relevant because it provides the conceptual basis for using pretrained backbones and domain-adaptation strategies when task-specific labels are limited. In the present thesis, these ideas are not used as generic performance tricks, but as part of a constrained design space in which model capacity, data availability, and supervision level must be balanced carefully.


\section{Research Gaps and Positioning of This Thesis}

The literature reviewed above suggests that the present thesis sits at the intersection of two partially connected research areas. On one side, CT-based mucus-plug studies have shown that patient-level mucus burden can be scored visually and that this burden carries clinically relevant information \cite{dunican2018mucus,dunican2018autopsy}. More recent work has extended this area in two directions: toward contemporary clinical and pathophysiologic study of mucus plugs \cite{hsieh2025smallairwaymucus}, and toward automated CT-based quantification, validation, detection, and segmentation \cite{vanderveer2025aicopdmucus,raut2025pragmaai,pu2026simulation}. On the other side, weakly supervised and multiple-instance CT learning studies have developed computational strategies for learning from coarse case-level labels and pooled local evidence \cite{qi2021drmil,xue2023copdmil,almakky2025weakly}. What remains less clearly addressed is the combination of these two perspectives: computational prediction of patient-level mucus burden from chest CT under weak supervision and limited data.

This gap is meaningful because the two literatures solve different parts of the problem. The CT mucus-plug literature establishes that the target is clinically interpretable, but it does not address how that target should be estimated automatically from CT. More recent automated mucus-plug studies move closer to computation, but they are still primarily concerned with quantification, validation, detection, or segmentation rather than weakly supervised patient-level regression \cite{vanderveer2025aicopdmucus,raut2025pragmaai,pu2026simulation}. By contrast, the weakly supervised CT literature offers relevant methodological patterns for learning from volumetric image evidence under coarse supervision, but it is mostly focused on case-level classification rather than regression of a clinically defined burden score \cite{qi2021drmil,xue2023copdmil,almakky2025weakly}. Existing work therefore does not directly address the setting studied in this thesis: patient-level mucus score prediction from chest CT when only coarse patient-level labels are available.

A second gap concerns the combination of supervision constraints and data scale. Many related medical-imaging studies either address different targets or operate in settings with annotation structures, data regimes, or model objectives that do not match patient-level mucus burden scoring. The present thesis instead studies a constrained setting in which the target is clinically meaningful, the supervision is weak, and the available cohort is small. This makes the problem distinct from both localized detection or segmentation tasks and more standard case-level classification settings.

The present thesis is therefore positioned as a machine-learning systems study built around an existing clinically meaningful target. It does not introduce a new mucus-plug scoring system, and it does not attempt plug-level detection or segmentation. Instead, it asks how far a slice-based, weakly supervised CT pipeline can go when only patient-level mucus scores are available. This makes the contribution mainly methodological and empirical, the thesis evaluates a constrained prediction pipeline and analyzes which design choices appear useful under limited-data conditions.
